
\chapter{Glosario de términos}

\DTLnewdb{terminos}

\DTLnewrow{terminos}
\DTLnewdbentry{terminos}{termino}{Catfishing/Catphishing}
\DTLnewdbentry{terminos}{definicion}{Práctica engañosa en la que una persona crea una identidad falsa en línea para establecer relaciones con otras personas,
generalmente con fines de manipulación emocional, fraude o explotación. Es común en aplicaciones de citas y redes
sociales, donde los estafadores se hacen pasar por alguien más para ganarse la confianza de la víctima y, en muchos
casos, obtener beneficios económicos o personales.}

\DTLnewrow{terminos}
\DTLnewdbentry{terminos}{termino}{Zero-shot}
\DTLnewdbentry{terminos}{definicion}{Es una técnica de prompt engineering que se identifica por redactar el prompt con una instrucción directa sin incluir ejemplos sobre la tarea asignada}

\DTLnewrow{terminos}
\DTLnewdbentry{terminos}{termino}{ventana de contexto}
\DTLnewdbentry{terminos}{definicion}{La ventana de contexto de un LLM se puede considerar como el equivalente de su memoria de trabajo. Determina la cantidad de conversación que puede llevar a cabo sin olvidar los detalles de anteriores en el intercambio. También determina el tamaño máximo de tokens, documentos o muestras de código que puede procesar.
\cite{ibmVentanaContexto}
\cite{AnthropicVentanaContexto}
}

\DTLnewrow{terminos}
\DTLnewdbentry{terminos}{termino}{Few - shot}
\DTLnewdbentry{terminos}{definicion}{Es una técnica de prompt engineering que es identificada por incluir ejemplos sobre la tarea que se le indica al modelo dentro del mismo prompt}

\DTLnewrow{terminos}
\DTLnewdbentry{terminos}{termino}{Lenguaje humano}
\DTLnewdbentry{terminos}{definicion}{istemas de signos que sirven para establecer comunicación entre personas, especialmente mediante la palabra hablada o escrita.}

\DTLnewrow{terminos}
\DTLnewdbentry{terminos}{termino}{Ciberfraude}
\DTLnewdbentry{terminos}{definicion}{Delito que implica el uso de tecnologías de la información y comunicación para cometer actos fraudulentos, como el robo de identidad, estafas financieras, phishing, entre otros. El ciberfraude puede manifestarse en diversas formas, incluyendo correos electrónicos engañosos, sitios web falsos y suplantación de identidad en redes sociales.}

\DTLnewrow{terminos}
\DTLnewdbentry{terminos}{termino}{Prompt}
\DTLnewdbentry{terminos}{definicion}{Instrucción o entrada de texto proporcionada a un modelo largo de lenguaje para generar una respuesta específica. La calidad y claridad del prompt influyen directamente en la precisión y relevancia de la salida generada por el modelo de inteligencia artificial.}

\DTLnewrow{terminos}
\DTLnewdbentry{terminos}{termino}{Prompt Engineering}
\DTLnewdbentry{terminos}{definicion}{Disciplina que se enfoca en diseñar y optimizar prompts para guiar eficazmente a modelos de lenguaje de gran escala (LLM) en la generación de respuestas deseadas. Implica comprender cómo formular instrucciones precisas que alineen las salidas del modelo con los objetivos específicos del usuario o aplicación. \cite{PromptEngByIBM}}

\DTLnewrow{terminos}
\DTLnewdbentry{terminos}{termino}{Fine Tuning}
\DTLnewdbentry{terminos}{definicion}{Proceso de adaptar un modelo de lenguaje previamente entrenado para tareas o dominios específicos mediante un entrenamiento adicional con datos relevantes. En el contexto de este proyecto, se destaca el uso de técnicas de prompt engineering para afinar el modelo sin necesidad de reentrenarlo completamente, lo que permite una personalización eficiente y menos costosa en términos computacionales. \cite{FineTuningByIbm}}

\DTLnewrow{terminos}
\DTLnewdbentry{terminos}{termino}{Large language model LLM}
\DTLnewdbentry{terminos}{definicion}{Son modelos de inteligencia artificial que suelen ser entrenados con grandes cantidades de información principalmente de textos con el fin de entender y manipular el lenguaje humano.}

\DTLnewrow{terminos}
\DTLnewdbentry{terminos}{termino}{Refinamiento de un LLM}
\DTLnewdbentry{terminos}{definicion}{Se refiere a el procedimiento de adaptar o mejorar el desempeño de un modelo previamente entrenado en tareas o dominios especificas, no confundir con el preentrenamiento de un LLM, proceso diferente donde se busca el aprendizaje de las bases del lenguaje por parte del modelo mediante una gran cantidad de datos}

\DTLsort{termino}{terminos}

\begin{enumerate}
\DTLforeach{terminos}{\labelItem=termino,\descItem=definicion}{
  \item \textbf{\labelItem}: \descItem
}
\end{enumerate}

Se invita a consultar la sección “marco teorico” para profundizar mas en cada uno de los terminos
